\documentclass{article}
\usepackage[utf8]{inputenc}
\usepackage{geometry}
 \geometry{
 a4paper,
 total={170mm,257mm},
 left=20mm,
 top=20mm,
 }
 \usepackage{graphicx}
 \usepackage{titling}
 \usepackage{booktabs}
 \usepackage{listings}
 \usepackage{xcolor}
 \usepackage{hyperref}
 \usepackage{fancyvrb}
 \usepackage{float}
 \usepackage{amsmath}
 \definecolor{darkblue}{rgb}{0.0, 0.0, 0.5}

\RecustomVerbatimEnvironment{verbatim}{Verbatim}{formatcom=\color{darkblue}}

\setlength{\parindent}{0pt} %lui dau dong


\title{\textbf{DEEP LEARNING} \\ 
{\Large \textbf{Labwork 4 - Neural Network}}}

\date{May 2026}
 
 \usepackage{fancyhdr}
\fancypagestyle{plain}{
    \fancyhf{} 
    \fancyfoot[L]{\thedate}
    \fancyfoot[R]{\thepage}
    \fancyhead[L]{University of Science and Technology of Hanoi}
    \fancyhead[R]{Deep Learning}
}
\makeatletter
\def\@maketitle{%
  \newpage
  % \null
  \vskip 1em %
  \begin{center}%
  \let \footnote \thanks
    {\LARGE \@title \par}%
    \vskip 1em%
    %{\large \@date}%
  \end{center}%
  \par
  \vskip 1em}
\makeatother

\usepackage{cmbright}

% Configuration for code blocks
\definecolor{codegreen}{rgb}{0,0.6,0}
\definecolor{backcolour}{rgb}{0.95,0.95,0.92}
\lstset{
    backgroundcolor=\color{backcolour},   
    commentstyle=\color{codegreen},
    basicstyle=\ttfamily\footnotesize,
    breaklines=true,
    showstringspaces=false
}

\begin{document}

\pagestyle{plain}
\maketitle

\noindent\begin{tabular}{@{}ll}
    Student: & Foo Cheung Nicolas - ES2540023 \\
\end{tabular}

\section{Algorithm Implementation}

The objective of this labwork is to implement a neural network from scratch using object-oriented programming. The neural network supports feedforward propagation and can solve the XOR logical problem.



\subsection{Network Architecture}

The neural network is composed of several layers. Each layer contains multiple neurons, and each neuron contains:

\begin{itemize}
    \item A list of weights
    \item A bias
    \item An activation function
\end{itemize}

The structure of the network is stored in a text file. Example:

\begin{verbatim}
3
2
2
1
\end{verbatim}

This represents:

\begin{itemize}
    \item 3 layers in total (input, hidden, output)
    \item 2 neurons in the input layer
    \item 2 neurons in the hidden layer
    \item 1 neuron in the output layer
\end{itemize}


\subsection{Feedforward Implementation}
For each neuron, the weighted sum is computed using:

\[
z = \sum_{i=1}^{n} w_i x_i + b
\]

where:

\begin{itemize}
    \item $w_i$ are the weights
    \item $x_i$ are the inputs
    \item $b$ is the bias
\end{itemize}

The activation output is then computed.

Initially, the sigmoid activation function was used:

\[
\sigma(z) = \frac{1}{1 + e^{-z}}
\]

However, the XOR weights provided in the slides were designed for a perceptron using a step activation function. Therefore, the step function was used for the final XOR experiment:

\[
f(z) =
\begin{cases}
1 & \text{if } z \geq 0 \\
0 & \text{otherwise}
\end{cases}
\]

\section{Object-Oriented Design}
The network architecture is read from a text file. The implementation contains three main classes:

\begin{itemize}
    \item Neuron
    \item Layer
    \item NeuralNetwork
\end{itemize}

\subsection{Neuron Class}

The \texttt{Neuron} class stores weights and bias, and computes the activation output.

\begin{lstlisting}
class Neuron:
    def __init__(self, weights, bias):
        self.weights = weights
        self.bias = bias
\end{lstlisting}

\subsection{Layer Class}

The \texttt{Layer} class contains multiple neurons and computes outputs for the entire layer.

\begin{lstlisting}
class Layer:
    def __init__(self, neurons):
        self.neurons = neurons
\end{lstlisting}


\subsection{NeuralNetwork Class}

The \texttt{NeuralNetwork} class contains all layers and performs feedforward propagation.

\begin{lstlisting}
class NeuralNetwork:
    def __init__(self, layers):
        self.layers = layers
\end{lstlisting}


\section{XOR Experiment}
The XOR problem was tested using the following weights and biases.

\subsection{Hidden Layer}

Neuron 1 (NAND):

\[
w = [-1, -1], \quad b = 1.5
\]

Neuron 2 (OR):

\[
w = [1, 1], \quad b = -0.5
\]

\subsection{Output Layer}

Output neuron (AND):

\[
w = [1, 1], \quad b = -1.5
\]

\section{Results}

The network was tested using the four XOR inputs.

\begin{center}
\begin{tabular}{|c|c|}
\hline
Input & Output \\
\hline
(0,0) & 0 \\
(0,1) & 1 \\
(1,0) & 1 \\
(1,1) & 0 \\
\hline
\end{tabular}
\end{center}

The neural network successfully reproduced the XOR logical function.


\end{document}